% =====================================================================
%  CSE 200 Online-1 (B1)  ---  "Laplace Analysis and Applications"
%  Faithful reproduction. Compile with pdflatex (run TWICE so that
%  \ref, \label, table of floats and the bibliography resolve).
% =====================================================================
\documentclass[10pt]{article}

% ---------- Page geometry ----------
\usepackage[a4paper,margin=1in]{geometry}

% ---------- Encoding & fonts ----------
\usepackage[T1]{fontenc}      % proper hyphenation + accented glyphs (Acad\'emie)
\usepackage[utf8]{inputenc}   % lets you type accented chars directly too

% ---------- Mathematics ----------
\usepackage{amsmath}          % align, gather, \dfrac, \text, \xrightarrow, cases
\usepackage{amssymb}          % \triangleright, \oplus, extra relation symbols
\usepackage{amsfonts}         % \mathcal{L} script letters etc.

% ---------- Graphics / figures ----------
\usepackage{graphicx}         % \includegraphics, \scalebox
\usepackage{wrapfig}          % wrapfigure  (text flows around the figure)
\usepackage{rotating}         % \rotatebox / sidewaysfigure for rotated floats
\usepackage[siunitx]{circuitikz} % draws the RC circuit natively

% ---------- Tables ----------
\usepackage{array}            % better column types, \arraystretch
\usepackage{multirow}         % \multirow spanning rows

% ---------- Colour & code ----------
\usepackage[dvipsnames]{xcolor} % named colours for text & listings
\usepackage{listings}         % Python source listing with syntax highlight

% ---------- Lists ----------
\usepackage{enumitem}         % custom list labels ($\oplus$, $\triangleright$)

% ---------- Authors with numbered affiliations ----------
\usepackage{authblk}          % \author[n]{...} + \affil[n]{...}

% ---------- Hyperlinks (load hyperref LAST) ----------
\usepackage{hyperref}
\hypersetup{
  colorlinks = true,
  linkcolor  = black,   % internal \ref links (Eq/Fig/Table)
  citecolor  = black,   % \cite links
  urlcolor   = blue     % external \href links
}

% ---------- Listings style for Python ----------
\lstdefinestyle{py}{
  language        = Python,
  basicstyle      = \ttfamily\small,
  keywordstyle    = \color{blue}\bfseries,
  commentstyle    = \color{OliveGreen}\itshape,
  stringstyle     = \color{BrickRed},
  numbers         = left,
  numberstyle     = \tiny\color{gray},
  numbersep       = 10pt,
  showstringspaces= false,
  frame           = single,
  breaklines      = true,
  captionpos      = b
}

% =====================================================================
\title{Laplace Analysis and Applications}
\author[1]{Pierre-Simon Laplace}
\author[2]{Oliver Heaviside}
\author[3]{Gustav Robert Kirchhoff}
\affil[1]{Acad\'emie des Sciences, Paris, France}
\affil[2]{Royal Society, London, United Kingdom}
\affil[3]{University of Heidelberg, Heidelberg, Germany}
\date{}                       % empty {} removes the date line
% =====================================================================

\begin{document}
\maketitle

% =====================================================================
\section{From Functions to Transforms}

\subsection{The Laplace Transform}
Let $f(t)$ be a function defined for $t \ge 0$. Its Laplace transform is
defined by
\begin{equation}
  \mathcal{L}\{f(t)\} = F(s) = \int_{0}^{\infty} e^{-st} f(t)\, dt,
  \qquad s > 0 .
  \label{eq:def}
\end{equation}
The transform changes a function of time into a function of the new
variable $s$.\footnote{The variable $s$ is commonly interpreted as a
complex-frequency variable.} Symbolically,
\[
  f(t) \xrightarrow{\;\mathcal{L}\;} F(s).
\]
The inverse Laplace transform reverses this process:
\begin{equation}
  f(t) = \mathcal{L}^{-1}\{F(s)\}.
  \label{eq:inv}
\end{equation}
The main purpose of the method is to replace differentiation and
integration in the time domain by simpler algebraic operations in the
transform domain \cite{kreyszig}. A concise overview is available from the
\href{https://mathworld.wolfram.com/LaplaceTransform.html}{Laplace
Transform entry on MathWorld}.

\subsection{A Simple Example}
For the constant function $f(t)=1$,
\begin{equation}
  \mathcal{L}\{1\}
    = \int_{0}^{\infty} e^{-st}\, dt
    = \left[ -\frac{1}{s} e^{-st} \right]_{0}^{\infty}
    = \frac{1}{s}.
  \label{eq:one}
\end{equation}
Similarly,
\begin{equation}
  \mathcal{L}\{t\} = \frac{1}{s^{2}}, \qquad
  \mathcal{L}\{t^{2}\} = \frac{2}{s^{3}}.
  \label{eq:powers}
\end{equation}

% =====================================================================
\section{Common Transform Equations}
Table~\ref{tab:transforms} lists several frequently used
Laplace-transform equations.

\begin{table}[htbp]
  \centering
  \caption{Selected Laplace-transform equations.}
  \label{tab:transforms}
  \renewcommand{\arraystretch}{1.6}   % taller rows so fractions breathe
  \begin{tabular}{|c|c|c|c|}
    \hline
    \multirow{2}{*}{\textbf{Family}}
      & \multicolumn{2}{c|}{\textbf{Transform Equation}}
      & \multirow{2}{*}{\textbf{Condition}} \\
    \cline{2-3}
      & \textbf{Time function} & \textbf{Laplace transform} & \\
    \hline
    \multirow{3}{*}{Algebraic}
      & $1$   & $\dfrac{1}{s}$        & $s > 0$ \\ \cline{2-4}
      & $t$   & $\dfrac{1}{s^{2}}$    & $s > 0$ \\ \cline{2-4}
      & $t^{n}$ & $\dfrac{n!}{s^{n+1}}$ & $n = 0,1,2,\dots$ \\
    \hline
    \multirow{2}{*}{Exponential}
      & $e^{at}$  & $\dfrac{1}{s-a}$ & $s > a$  \\ \cline{2-4}
      & $e^{-at}$ & $\dfrac{1}{s+a}$ & $s > -a$ \\
    \hline
    \multirow{2}{*}{Oscillatory}
      & $\sin(\omega t)$ & $\dfrac{\omega}{s^{2}+\omega^{2}}$ & $\omega > 0$ \\ \cline{2-4}
      & $\cos(\omega t)$ & $\dfrac{s}{s^{2}+\omega^{2}}$      & $\omega > 0$ \\
    \hline
    \multicolumn{2}{|c|}{\textbf{Unit-step function }$u(t-a)$}
      & $\dfrac{e^{-as}}{s}$ & $a > 0$ \\
    \hline
  \end{tabular}
\end{table}

The formulas in Table~\ref{tab:transforms} allow many functions to be
transformed without repeatedly evaluating improper integrals.

% =====================================================================
\section{Application to an RC Circuit}

\subsection*{Circuit Model}

\begin{wrapfigure}{r}{0.42\textwidth}
  \centering
  \begin{circuitikz}[scale=0.85, transform shape]
    \draw (0,0)
      to[V, l=$v_{in}(t)$, invert] (0,3)     % source, + on top
      to[short, i=$i(t)$] (1.2,3)
      to[R, l=$R$] (4,3)                      % resistor on top branch
      -- (5.2,3)
      to[C, l_=$C$, v=$v_C(t)$] (5.2,0)       % capacitor on right branch
      -- (0,0);                               % close the loop
  \end{circuitikz}
  \caption{A simple series RC circuit.}
  \label{fig:rc}
\end{wrapfigure}

Figure~\ref{fig:rc} shows a simple \texttt{resistor-{}-capacitor circuit}
connected to an input voltage source. The resistor limits the current,
while the capacitor stores electrical energy.

For a resistor $R$ and capacitor $C$, Kirchhoff's voltage law gives
\begin{equation}
  RC\,\frac{dv_{C}(t)}{dt} + v_{C}(t) = v_{\text{in}}(t),
  \label{eq:kvl}
\end{equation}
where $v_{C}(t)$ is the voltage across the capacitor.

Assuming $v_{C}(0)=0$, taking the Laplace transform of
Equation~\eqref{eq:kvl} gives
\begin{align}
  RC\bigl[sV_{C}(s) - v_{C}(0)\bigr] + V_{C}(s) &= V_{\text{in}}(s), \nonumber\\
  \Rightarrow RCsV_{C}(s) + V_{C}(s)            &= V_{\text{in}}(s), \nonumber\\
  \Rightarrow V_{C}(s)\,(RCs + 1)               &= V_{\text{in}}(s).
  \label{eq:algebra}
\end{align}

The transfer function is therefore
\begin{equation}
  H(s) = \frac{V_{C}(s)}{V_{\text{in}}(s)} = \frac{1}{RCs + 1}.
  \label{eq:tf}
\end{equation}
For a constant input $V_{0}$,
\[
  V_{\text{in}}(s) = \frac{V_{0}}{s}.
\]
Thus,
\begin{equation}
  V_{C}(s) = \frac{V_{0}}{s(RCs + 1)}.
  \label{eq:vcs}
\end{equation}
Taking the inverse transform gives
\begin{equation}
  v_{C}(t) = V_{0}\left( 1 - e^{-t/(RC)} \right).
  \label{eq:vct}
\end{equation}
The quantity $\tau = RC$ is called the time constant of the circuit.
A larger value of $\tau$ causes the capacitor to charge more
slowly \cite{ogata}.

% =====================================================================
\section{How Laplace Analysis Is Used}
The Laplace transform is useful in many areas:

\begin{description}[leftmargin=0pt, style=nextline]
  \item[Differential equations]
    \begin{itemize}[label=$\oplus$]
      \item initial conditions are included automatically;
      \item derivatives become algebraic expressions.
    \end{itemize}
  \item[Electrical circuits]
    \begin{itemize}[label=$\oplus$]
      \item voltages and currents are studied in the $s$-domain;
      \item components lead to transfer functions.
    \end{itemize}
  \item[Control and signal systems]
    \begin{itemize}[label=$\oplus$]
      \item stability can be studied through poles;
      \item common test inputs include:
        \begin{itemize}[label=$\triangleright$]
          \item the unit-step input;
          \item the impulse input;
          \item sinusoidal inputs.
        \end{itemize}
    \end{itemize}
\end{description}

In the time domain, a system may be described by a differential equation,
whereas in the Laplace domain, the same system is often described by a
rational algebraic expression.

% =====================================================================
\section{A Short Python Implementation}
Listing~\ref{lst:rc} evaluates the capacitor voltage at a finite set of
time instants using an external Python file.

\begin{lstlisting}[style=py, caption={Computing samples of the RC step response.}, label={lst:rc}]
from math import exp


def rc_response(resistance, capacitance, input_voltage, times):
    """Return capacitor-voltage samples for an RC step input."""
    if resistance <= 0 or capacitance <= 0:
        raise ValueError("R and C must be positive")

    tau = resistance * capacitance
    values = []

    for time in times:
        voltage = input_voltage * (1 - exp(-time / tau))
        values.append(voltage)

    return values


times = [0.0, 0.5, 1.0, 2.0]
print(rc_response(2.0, 0.5, 5.0, times))
\end{lstlisting}

Equation~\eqref{eq:vct} and Listing~\ref{lst:rc} describe the same
computation.

% =====================================================================
\begin{thebibliography}{9}
  \bibitem{kreyszig} E.~Kreyszig, \textit{Advanced Engineering
    Mathematics}, 10th ed. John Wiley \& Sons, 2011.
  \bibitem{ogata} K.~Ogata, \textit{Modern Control Engineering},
    5th ed. Prentice Hall, 2010.
\end{thebibliography}

\end{document}
